\subsection{SIMPLE ELEMENTS}

DEFINITION 2. An element \(h\) of \(\mathfrak{g}\) is said to be simple if there exist \(x, y \in \mathfrak{g}\) such that \((x, h, y)\) is an \(\mathfrak{sl}_2\) -triplet in \(\mathfrak{g}\) .

We also say that \(h\) is the simple element of the \(\mathfrak{sl}_2\) -triplet \((x, h, y)\) .

PROPOSITION 4. Let \(h\) be a non-zero element of \(\mathfrak{g}\) . Then \(h\) is simple if and only if there exists \(x \in \mathfrak{g}\) such that \([h, x] = 2x\) and \(h \in (\operatorname{ad} x)(\mathfrak{g})\) .

The condition is clearly necessary. It is sufficient by Lemma 6.

PROPOSITION 5. Assume that g is splittable semi-simple. Let h be a splitting Cartan subalgebra of g, R the set of roots of \((\mathfrak{g},\mathfrak{h})\) , and B a basis of R. Let h be a simple element of g belonging to h. Then h is conjugate under \(\operatorname{Aut}_{e}(\mathfrak{g},\mathfrak{h})\) to an element \(h'\) of h such that \(\alpha(h') \in \{0,1,2\}\) for all \(\alpha \in B\) .

The eigenvalues of \(\mathrm{ad}_{\mathfrak{g}}h\) belong to \(\mathbf{Z}\) (§1, no. 2, Cor. of Prop. 2). Hence \(h \in \mathfrak{h}_{\mathbf{Q}}\) . There exists an element \(w\) of the Weyl group of \((\mathfrak{g}, \mathfrak{h})\) such that \(\alpha(wh) \geq 0\) for all \(\alpha \in \mathrm{B}\) (Chap. VI, §1, no. 5, Th. 2 (i)). In view of §2, no. 2, Cor. of Th. 2, we are reduced to the case in which \(\alpha(h) \in \mathbf{N}\) for all \(\alpha \in \mathrm{B}\) . Let \(\mathrm{R}_{+}\) be the set of positive roots relative to B, and \(\mathrm{R}_{-} = -\mathrm{R}_{+}\) . There exists an \(\mathfrak{sl}_2\) -triplet in \(\mathfrak{g}\) of the form \((x, h, y)\) . Let T be the set of roots \(\beta\) such that \(\beta(h) = -2\) . Then \(\mathrm{T} \subset \mathrm{R}_{-}\) and \(y \in \sum_{\beta \in \mathrm{T}} \mathfrak{g}^{\beta}\) . Assume that there exists \(\alpha \in \mathrm{B}\) such that \(\alpha(h) > 2\) . For all \(\beta \in \mathrm{T}\) , we have \((\alpha + \beta)(h) > 0\) , so \(\alpha + \beta \notin \mathrm{R}_{-}\) and \(\alpha + \beta \neq 0\) ; on the other hand, since \(\beta \in \mathrm{R}_{-}\) and \(\alpha \in \mathrm{B}\) , we have \(\alpha + \beta \notin \mathrm{R}_{+}\) ; hence \(\alpha + \beta \notin \mathrm{R} \cup \{0\}\) , so \([\mathfrak{g}^{\alpha}, \mathfrak{g}^{\beta}] = 0\) . Thus, \([y, \mathfrak{g}^{\alpha}] = 0\) . But \(\mathrm{ad}_{\mathfrak{g}}y|\mathfrak{g}^{\alpha}\) is injective since \(\alpha(h) > 0\) (§1, no. 2, Cor. of Prop. 2). This contradiction proves that \(\alpha(h) \leq 2\) for all \(\alpha \in \mathrm{B}\) .

COROLLARY. If \(k\) is algebraically closed and if \(\mathfrak{g}\) is semi-simple of rank \(l\) , the number of conjugacy classes of simple elements of \(\mathfrak{g}\) , relative to \(\operatorname{Aut}_e(\mathfrak{g})\) , is at most \(3^l\) .

Indeed, every semi-simple element of \(\mathfrak{g}\) is conjugate under \(\mathrm{Aut}_e(\mathfrak{g})\) to an element of \(\mathfrak{h}\) .

Lemma 8. Assume that \(k\) is algebraically closed and that \(\mathfrak{g}\) is semi-simple. Let \(h\) be a semi-simple element of \(\mathfrak{g}\) such that the eigenvalues of \(\operatorname{ad} h\) are rational. Let \(\mathfrak{g}^0 = \operatorname{Ker}(\operatorname{ad} h)\) , \(\mathfrak{g}^2 = \operatorname{Ker}(\operatorname{ad} h - 2)\) . Let \(G_h\) be the set of elementary automorphisms of \(\mathfrak{g}\) leaving \(h\) fixed. Let \(x \in \mathfrak{g}^2\) be such that \([x, \mathfrak{g}^0] = \mathfrak{g}^2\) . Then \(G_h x\) contains a subset of \(\mathfrak{g}^2\) that is dense and open in the Zariski topology.

Let \(\mathfrak{h}\) be a Cartan subalgebra of \(\mathfrak{g}^0\) . This is a Cartan subalgebra of \(\mathfrak{g}\) containing \(h\) (Chap. VII, §2, no. 3, Prop. 10). We have \(h \in \mathfrak{h}_{\mathbf{Q}}\) . Let R be the root system of \((\mathfrak{g}, \mathfrak{h})\) , Q the group of radical weights. There exists a basis B of R such that \(\alpha(h) \geq 0\) for all \(\alpha \in B\) .

Let \(\mathrm{U}\) be the set of \(z \in \mathfrak{h}\) such that \(\alpha(z) \neq 0\) for all \(\alpha \in \mathrm{B}\) . Let \((H_{\alpha}')_{\alpha \in \mathrm{B}}\) be the basis of \(\mathfrak{h}\) dual to \(\mathrm{B}\) . If \(z \in \mathrm{U}\) , there exists a homomorphism from \(\mathrm{Q}\) to \(k^*\) that takes any \(\gamma \in \mathrm{Q}\) to \(\prod_{\alpha \in \mathrm{B}} \alpha(z)^{\gamma(H_{\alpha}')}\) . By §5, Prop. 2 and 4, the endomorphism \(\varphi(z)\) of the vector space \(\mathfrak{g}\) which induces on \(\mathfrak{g}^{\gamma}\) the homothety with ratio \(\prod_{\alpha \in \mathrm{B}} \alpha(z)^{\gamma(H_{\alpha}')}\) is an elementary automorphism of \(\mathfrak{g}\) , which clearly belongs to \(G_h\) .

Let \(s \in \mathfrak{h}\) . If \(\gamma \in \mathrm{R}\) is such that \(\mathfrak{g}^{\gamma} \cap \mathfrak{g}^{2} \neq 0\) ,

\[
2 = \gamma (h) = \gamma \left(\sum_ {\alpha \in \mathrm{B}} \alpha (h) H _ {\alpha} ^ {\prime}\right) = \sum_ {\alpha \in \mathrm{B}} \alpha (h) \gamma (H _ {\alpha} ^ {\prime});
\]

since \(\alpha(h) \geq 0\) for all \(\alpha \in \mathrm{B}\) , and since the \(\gamma(H_{\alpha}^{\prime})\) are integers either all \(\geq 0\) or all \(\leq 0\) , we have \(\gamma(H_{\alpha}^{\prime}) \in \mathbf{N}\) for all \(\alpha \in \mathrm{B}\) . Thus, we can consider (for \(z \in \mathfrak{h}\) ) the endomorphism \(\psi(z)\) of the vector space \(\mathfrak{g}^2\) that induces on \(\mathfrak{g}^{\gamma} \cap \mathfrak{g}^2\) the homothety with ratio \(\prod_{\alpha \in \mathrm{B}} \alpha(z)^{\gamma(H_{\alpha}^{\prime})}\) . The map \(z \mapsto \psi(z)\) from \(\mathfrak{h}\) to \(\operatorname{End}(\mathfrak{g}^2)\) is polynomial. For \(z \in \mathrm{U}\) , we have \(\psi(z) = \varphi(z)|\mathfrak{g}^2\) .

Let \(\gamma_1, \ldots, \gamma_r\) be the distinct roots of \((\mathfrak{g}, \mathfrak{h})\) vanishing on \(h\) . If \(y_1 \in \mathfrak{g}^{\gamma_1}, \ldots, y_r \in \mathfrak{g}^{\gamma_r}\) , we have \(e^{\mathrm{ad} y_1} \ldots e^{\mathrm{ad} y_r} \in \mathrm{G}_h\) . We can thus define a map \(\rho\) from \(\mathfrak{h} \times \mathfrak{g}^{\gamma_1} \times \cdots \times \mathfrak{g}^{\gamma_r}\) to \(\mathfrak{g}^2\) by putting

\[
\rho (z, y _ {1}, \dots , y _ {r}) = \psi (z) e ^ {\operatorname{ad} y _ {1}} \dots e ^ {\operatorname{ad} y _ {r}} x
\]

for \(z \in \mathfrak{h}\) , \(y_1 \in \mathfrak{g}^{\gamma_1}, \ldots, y_r \in \mathfrak{g}^{\gamma_r}\) . This map is polynomial, and \(\rho(\mathrm{U}, \mathfrak{g}^{\gamma_1}, \ldots, \mathfrak{g}^{\gamma_r}) \subset \mathrm{G}_h x\) . By Chap. VII, App. I, Prop. 3 and 4, it suffices to prove that the tangent linear map of \(\rho\) is surjective at some point.

Now let T be the tangent linear map of \(z \mapsto \psi(z)\) at \(h_{0} = \sum_{\alpha \in B} H_{\alpha}'\) . Then \(\mathrm{T}(z)\) is the endomorphism of \(g^{2}\) that induces on \(g^{\gamma} \cap g^{2}\) the homothety with ratio

\[
\sum_ {\alpha \in \mathrm{B}} \gamma (H _ {\alpha} ^ {\prime}) \alpha (h _ {0}) ^ {\gamma (H _ {\alpha} ^ {\prime}) - 1} \alpha (z) \prod_ {\beta \in \mathrm{B}, \beta \neq \alpha} \beta (h _ {0}) ^ {\gamma (H _ {\alpha} ^ {\prime})} = \sum_ {\alpha \in \mathrm{B}} \gamma (H _ {\alpha} ^ {\prime}) \alpha (z) = \gamma (z).
\]

Thus, the tangent linear map of \(z \mapsto \rho(z,0,\ldots,0)\) at \(h_{0}\) is the map \(z \mapsto [z,x]\) ; its image is \([x,h]\) . The tangent linear map at 0 of the map \(y_{1} \mapsto \rho(h_{0},y_{1},0,\ldots,0)\) is the map \(y_{1} \mapsto \psi(h_{0})[y_{1},x]\) ; this last map has image \(\psi(h_{0})[x,\mathfrak{g}^{\gamma_{1}}] = [x,\mathfrak{g}^{\gamma_{1}}]\) . Similarly, the tangent linear map at 0 of the map \(y_{i} \mapsto \rho(h_{0},0,\ldots,0,y_{i},0,\ldots,0)\) has image \([x,\mathfrak{g}^{\gamma_{i}}]\) . Finally, the tangent linear map of \(\rho\) at \((h_{0},0,\ldots,0)\) has image

\[
[ x, \mathfrak {h} + \mathfrak {g} ^ {\gamma_ {1}} + \dots + \mathfrak {g} ^ {\gamma_ {r}} ] = [ x, \mathfrak {g} ^ {0} ] = \mathfrak {g} ^ {2}.\tag{Q.E.D.}
\]

*The group \(\mathrm{G}_h\) is an algebraic group with Lie algebra ad \(\mathfrak{g}^0_*\)

PROPOSITION 6. Assume that \(k\) is algebraically closed and that \(\mathfrak{g}\) is semi-simple. Let \(G\) be a group of automorphisms of \(\mathfrak{g}\) containing \(\operatorname{Aut}_e(\mathfrak{g})\) . Let \((x, h, y)\) and \((x', h', y')\) be \(\mathfrak{sl}_2\) -triplets in \(\mathfrak{g}\) . The following conditions are equivalent:

\begin{enumerate}
\def\labelenumi{(\roman{enumi})}
\item
  h and \(h'\) are G-conjugate;
\item
  \((x,h,y)\) and \((x^{\prime},h^{\prime},y^{\prime})\) are G-conjugate.
\end{enumerate}

We only have to prove the implication (i) \(\Longrightarrow\) (ii), and we are reduced immediately to the case in which \(h = h'\) . Introduce \(g^{2}\) and \(G_{h}\) as in Lemma 8. We have \(x \in g^{2}\) , and \([x, g^{0}] = g^{2}\) by §1, no. 2, Cor. of Prop. 2. Hence \(G_{h}x\) contains a subset of \(g^{2}\) that is dense and open in the Zariski topology, and so does \(G_{h}x'\) . So there exists \(a \in G_{h}\) such that \(a(x) = x'\) . We have \(a(h) = h\) , and consequently \(a(y) = y'\) (no. 1, Lemma 1).

COROLLARY 1. The map which associates to any \(\mathfrak{sl}_2\) -triplet its simple element defines by passage to the quotients a bijection from the set of G-conjugacy classes of \(\mathfrak{sl}_2\) -triplets to the set of G-conjugacy classes of simple elements.

COROLLARY 2. If \(\operatorname{rk}(\mathfrak{g}) = l\) , the number of conjugacy classes, relative to \(\operatorname{Aut}_e(\mathfrak{g})\) , of non-zero nilpotent elements of \(\mathfrak{g}\) is at most \(3^l\) .

This follows from Cor. 1, the Cor. of Prop. 2, and the Cor. of Prop. 5.

COROLLARY 3. If \(\mathrm{rk}(\mathfrak{g}) = l\) , the number of conjugacy classes, relative to \(\mathrm{Aut}_e(\mathfrak{g})\) , of subalgebras of \(\mathfrak{g}\) isomorphic to \(\mathfrak{sl}(2,k)\) is at most \(3^l\) .

This follows from Cor. 1, Prop. 1, and the Cor. of Prop. 5.

